% Options for packages loaded elsewhere
\PassOptionsToPackage{unicode}{hyperref}
\PassOptionsToPackage{hyphens}{url}
%
\documentclass[
  10pt,
]{article}
\usepackage{amsmath,amssymb}
\usepackage{iftex}
\ifPDFTeX
  \usepackage[T1]{fontenc}
  \usepackage[utf8]{inputenc}
  \usepackage{textcomp} % provide euro and other symbols
\else % if luatex or xetex
  \usepackage{unicode-math} % this also loads fontspec
  \defaultfontfeatures{Scale=MatchLowercase}
  \defaultfontfeatures[\rmfamily]{Ligatures=TeX,Scale=1}
\fi
\usepackage{lmodern}
\ifPDFTeX\else
  % xetex/luatex font selection
\fi
% Use upquote if available, for straight quotes in verbatim environments
\IfFileExists{upquote.sty}{\usepackage{upquote}}{}
\IfFileExists{microtype.sty}{% use microtype if available
  \usepackage[]{microtype}
  \UseMicrotypeSet[protrusion]{basicmath} % disable protrusion for tt fonts
}{}
\makeatletter
\@ifundefined{KOMAClassName}{% if non-KOMA class
  \IfFileExists{parskip.sty}{%
    \usepackage{parskip}
  }{% else
    \setlength{\parindent}{0pt}
    \setlength{\parskip}{6pt plus 2pt minus 1pt}}
}{% if KOMA class
  \KOMAoptions{parskip=half}}
\makeatother
\usepackage{xcolor}
\usepackage[margin=22mm]{geometry}
\usepackage{color}
\usepackage{fancyvrb}
\newcommand{\VerbBar}{|}
\newcommand{\VERB}{\Verb[commandchars=\\\{\}]}
\DefineVerbatimEnvironment{Highlighting}{Verbatim}{commandchars=\\\{\}}
% Add ',fontsize=\small' for more characters per line
\newenvironment{Shaded}{}{}
\newcommand{\AlertTok}[1]{\textcolor[rgb]{1.00,0.00,0.00}{\textbf{#1}}}
\newcommand{\AnnotationTok}[1]{\textcolor[rgb]{0.38,0.63,0.69}{\textbf{\textit{#1}}}}
\newcommand{\AttributeTok}[1]{\textcolor[rgb]{0.49,0.56,0.16}{#1}}
\newcommand{\BaseNTok}[1]{\textcolor[rgb]{0.25,0.63,0.44}{#1}}
\newcommand{\BuiltInTok}[1]{\textcolor[rgb]{0.00,0.50,0.00}{#1}}
\newcommand{\CharTok}[1]{\textcolor[rgb]{0.25,0.44,0.63}{#1}}
\newcommand{\CommentTok}[1]{\textcolor[rgb]{0.38,0.63,0.69}{\textit{#1}}}
\newcommand{\CommentVarTok}[1]{\textcolor[rgb]{0.38,0.63,0.69}{\textbf{\textit{#1}}}}
\newcommand{\ConstantTok}[1]{\textcolor[rgb]{0.53,0.00,0.00}{#1}}
\newcommand{\ControlFlowTok}[1]{\textcolor[rgb]{0.00,0.44,0.13}{\textbf{#1}}}
\newcommand{\DataTypeTok}[1]{\textcolor[rgb]{0.56,0.13,0.00}{#1}}
\newcommand{\DecValTok}[1]{\textcolor[rgb]{0.25,0.63,0.44}{#1}}
\newcommand{\DocumentationTok}[1]{\textcolor[rgb]{0.73,0.13,0.13}{\textit{#1}}}
\newcommand{\ErrorTok}[1]{\textcolor[rgb]{1.00,0.00,0.00}{\textbf{#1}}}
\newcommand{\ExtensionTok}[1]{#1}
\newcommand{\FloatTok}[1]{\textcolor[rgb]{0.25,0.63,0.44}{#1}}
\newcommand{\FunctionTok}[1]{\textcolor[rgb]{0.02,0.16,0.49}{#1}}
\newcommand{\ImportTok}[1]{\textcolor[rgb]{0.00,0.50,0.00}{\textbf{#1}}}
\newcommand{\InformationTok}[1]{\textcolor[rgb]{0.38,0.63,0.69}{\textbf{\textit{#1}}}}
\newcommand{\KeywordTok}[1]{\textcolor[rgb]{0.00,0.44,0.13}{\textbf{#1}}}
\newcommand{\NormalTok}[1]{#1}
\newcommand{\OperatorTok}[1]{\textcolor[rgb]{0.40,0.40,0.40}{#1}}
\newcommand{\OtherTok}[1]{\textcolor[rgb]{0.00,0.44,0.13}{#1}}
\newcommand{\PreprocessorTok}[1]{\textcolor[rgb]{0.74,0.48,0.00}{#1}}
\newcommand{\RegionMarkerTok}[1]{#1}
\newcommand{\SpecialCharTok}[1]{\textcolor[rgb]{0.25,0.44,0.63}{#1}}
\newcommand{\SpecialStringTok}[1]{\textcolor[rgb]{0.73,0.40,0.53}{#1}}
\newcommand{\StringTok}[1]{\textcolor[rgb]{0.25,0.44,0.63}{#1}}
\newcommand{\VariableTok}[1]{\textcolor[rgb]{0.10,0.09,0.49}{#1}}
\newcommand{\VerbatimStringTok}[1]{\textcolor[rgb]{0.25,0.44,0.63}{#1}}
\newcommand{\WarningTok}[1]{\textcolor[rgb]{0.38,0.63,0.69}{\textbf{\textit{#1}}}}
\setlength{\emergencystretch}{3em} % prevent overfull lines
\providecommand{\tightlist}{%
  \setlength{\itemsep}{0pt}\setlength{\parskip}{0pt}}
\setcounter{secnumdepth}{-\maxdimen} % remove section numbering
\ifLuaTeX
  \usepackage{selnolig}  % disable illegal ligatures
\fi
\usepackage{bookmark}
\IfFileExists{xurl.sty}{\usepackage{xurl}}{} % add URL line breaks if available
\urlstyle{same}
\hypersetup{
  hidelinks,
  pdfcreator={LaTeX via pandoc}}

\author{}
\date{}

\begin{document}

\section{DDTA Decision phase - work
plan}\label{ddta-decision-phase---work-plan}

\textbf{Non-canonical research plan - checkpoint before facial-access
Decisions}

\subsection{Checkpoint purpose}\label{checkpoint-purpose}

This file freezes the intended experimental sequence before authoring
any Decision for the independent facial-access / controlled-access demo.

Expected local repository baseline for applying this checkpoint: commit
\texttt{2db7604}
(\texttt{research:\ add\ MR\ teachability\ and\ pre-close\ projection\ study}).
The full 40-character SHA was not captured in the prior transcript, so
application commands verify the short HEAD before extraction.

The current Decision construction candidate was built from the
ThreatForge MR-0001 + MR-0002 historical ADR corpus only. The
independent facial-access Decisions have not yet been written, and
ThreatForge MR-0003, MR-0004 and MR-0005 remain untouched sequential
holdouts.

\subsection{Current frozen construction-candidate
baseline}\label{current-frozen-construction-candidate-baseline}

The Decision-specific core is closed for the construction candidate:

\begin{Shaded}
\begin{Highlighting}[]
\NormalTok{Decision}
\NormalTok{|{-} title}
\NormalTok{|{-} context}
\NormalTok{|{-} decision}
\NormalTok{\textasciigrave{}{-} consequences}
\end{Highlighting}
\end{Shaded}

All four fields are semantically required and non-null. Parentage to one
Macro Requirement is an explicit semantic relation. Lifecycle/status and
serialization storage remain separate concerns.

The following construction conclusions are also frozen unless a concrete
independent counterexample reopens them:

\begin{enumerate}
\def\labelenumi{\arabic{enumi}.}
\tightlist
\item
  Context is decision-local rather than a repetition of the parent MR.
\item
  One Decision expresses one coherent chosen position.
\item
  Consequences expose material effects/trade-offs of the choice.
\item
  Non-goals is not a Decision-core patch field; misleading Context +
  Decision must be rewritten.
\item
  Separate mandatory \texttt{rationale} and embedded
  \texttt{alternatives} are not currently justified.
\item
  DDTA/metamodel, governed project, ThreatForge tool and ThreatForge
  project operations have distinct semantic ownership.
\item
  A ThreatForge ADR may decide how the tool supports general semantics,
  but must not silently define those general semantics.
\item
  Cross-document persistence problems are not solved with topic-specific
  Decision fields.
\end{enumerate}

Cross-cutting representation decisions are also frozen for this
construction candidate:

\begin{itemize}
\tightlist
\item
  complete canonical field shape;
\item
  explicit field-specific nullability / canonical empty values;
\item
  semantic/representation separation;
\item
  executable projections derived from one governed canonical document
  model.
\end{itemize}

\subsection{Experimental discipline}\label{experimental-discipline}

The next phases must preserve evidence independence.

\begin{itemize}
\tightlist
\item
  Do not use ThreatForge MR-0003/MR-0004/MR-0005 while authoring the
  facial-access Decisions.
\item
  Do not refactor or rewrite ThreatForge implementation or canonical
  project documents during this conceptual study.
\item
  Do not define Requirement semantics merely to make the Decision model
  easier to validate.
\item
  Do not introduce lifecycle/status mechanics during the Decision
  semantic test.
\item
  Do not add fields to Decision merely because one case needs persistent
  applicability, vocabulary, patterns, traceability or another
  cross-cutting concern.
\item
  Technical/domain vocabulary is allowed when intrinsic to the actual
  choice; premature metamodel vocabulary must not become project
  semantics.
\item
  Every change to the construction candidate after this checkpoint
  requires a concrete counterexample and an explicit regression over all
  previously used cases.
\end{itemize}

\subsection{Phase D1 - author independent facial-access
Decisions}\label{phase-d1---author-independent-facial-access-decisions}

Use the four already-studied Macro Requirements of the independent
facial-access / controlled-access example:

\begin{enumerate}
\def\labelenumi{\arabic{enumi}.}
\tightlist
\item
  MR-0001 - Accesso controllato all'area riservata
\item
  MR-0002 - Gestione delle persone autorizzate
\item
  MR-0003 - Riconoscimento facciale per la verifica della persona
\item
  MR-0004 - Uso responsabile dei dati biometrici e delle decisioni
  automatiche
\end{enumerate}

Author a deliberately small but varied Decision corpus, approximately
6-8 Decisions across the four MRs.

Authoring constraints:

\begin{itemize}
\tightlist
\item
  use only \texttt{Title}, \texttt{Context}, \texttt{Decision},
  \texttt{Consequences} for Decision-specific semantics;
\item
  write each field explicitly; no optional-by-omission representation;
\item
  avoid Requirement-level independently checkable detail unless that
  precision is intrinsic to the choice itself;
\item
  avoid premature architecture (camera topology, protocols, databases,
  thresholds, APIs, deployment) unless the Decision genuinely chooses
  it;
\item
  include at least one non-software / policy-or-method Decision so the
  model is not validated only on architecture choices;
\item
  keep the corpus independent from the historical ThreatForge wording.
\end{itemize}

Deliverable: a readable demo document containing the candidate
facial-access Decisions, without changing the Decision metamodel during
initial authoring.

\subsection{Phase D2 - analyze the facial-access Decisions against the
frozen
candidate}\label{phase-d2---analyze-the-facial-access-decisions-against-the-frozen-candidate}

Analyze each new Decision using the same tests:

\begin{enumerate}
\def\labelenumi{\arabic{enumi}.}
\tightlist
\item
  parent/MR narrowing;
\item
  decision-local Context;
\item
  explicit coherent choice;
\item
  consequence awareness;
\item
  semantic-owner correctness;
\item
  Decision-vs-downstream boundary;
\item
  no semantic patching;
\item
  model/tool independence;
\item
  survivability challenge without ad-hoc fields;
\item
  compatibility with complete-shape / derived executable representation.
\end{enumerate}

For every Decision produce one of two outcomes:

\begin{Shaded}
\begin{Highlighting}[]
\NormalTok{fits construction candidate}

\NormalTok{or}

\NormalTok{counterexample {-}\textgreater{} candidate correction required}
\end{Highlighting}
\end{Shaded}

A counterexample must state exactly which closed rule fails and why the
information cannot be represented without loss or distortion.

\subsection{Phase D3 - regression and Decision candidate
closure}\label{phase-d3---regression-and-decision-candidate-closure}

If D2 exposes a counterexample:

\begin{enumerate}
\def\labelenumi{\arabic{enumi}.}
\tightlist
\item
  revise only the minimum affected rule;
\item
  record the evidence that caused the revision;
\item
  rerun the affected facial-access cases;
\item
  rerun all 16 ThreatForge construction ADRs against the corrected
  model.
\end{enumerate}

If no unresolved counterexample remains, mark the Decision model as
\textbf{candidate closed after independent example validation}.

This does not yet establish universal validity; it authorizes the
sequential holdout phase.

\subsection{Phase D4 - ThreatForge MR-0003
holdout}\label{phase-d4---threatforge-mr-0003-holdout}

Read and analyze MR-0003 Decisions for the first time as holdout
evidence.

\begin{itemize}
\tightlist
\item
  Do not rewrite the model pre-emptively.
\item
  Classify every mismatch as historical ADR defect, ownership/conflation
  defect, downstream leakage, or genuine metamodel counterexample.
\item
  Modify the model only for the last category.
\item
  If modified, regress ThreatForge MR-0001/MR-0002 plus the
  facial-access corpus before continuing.
\end{itemize}

\subsection{Phase D5 - ThreatForge MR-0004
holdout}\label{phase-d5---threatforge-mr-0004-holdout}

Repeat the same procedure using MR-0004 only after MR-0003 is resolved.
Any model modification requires regression on all earlier evidence.

\subsection{Phase D6 - ThreatForge MR-0005
holdout}\label{phase-d6---threatforge-mr-0005-holdout}

Repeat the same procedure using MR-0005 only after MR-0004 is resolved.
This is the final Decision holdout in the planned sequence.

\subsection{Phase D7 - Decision-phase
synthesis}\label{phase-d7---decision-phase-synthesis}

After all three holdouts:

\begin{itemize}
\tightlist
\item
  state the final Decision semantic model and invariants supported by
  the evidence;
\item
  separate stable general rules from ThreatForge case-study lessons;
\item
  record remaining cross-document constraints that cannot be solved at
  Decision level;
\item
  preserve the complete-shape / semantic-representation /
  executable-projection decisions for reuse by later document models;
\item
  update authoring guidance and anti-patterns only after the final
  synthesis.
\end{itemize}

\subsection{Phase R1 - move to Requirement
semantics}\label{phase-r1---move-to-requirement-semantics}

Only after Decision-phase synthesis, begin the next-level Requirement
study.

Priority questions already carried forward:

\begin{itemize}
\tightlist
\item
  what makes an obligation independently checkable without reducing
  every Requirement to one implementation file;
\item
  how pervasive/cross-cutting obligations are represented and covered;
\item
  how persistent Decisions retain downstream effect without
  inheritance-by-copy;
\item
  how single-source governed knowledge and effective governed context
  are represented;
\item
  how humans, tools and LLMs recover applicable context without relying
  on memory;
\item
  how Requirement semantics map into complete canonical representations
  and executable projections.
\end{itemize}

\subsection{Deferred beyond the current
plan}\label{deferred-beyond-the-current-plan}

\begin{itemize}
\tightlist
\item
  lifecycle/status transitions and historical governance;
\item
  concrete YAML/Markdown serialization;
\item
  ThreatForge refactoring/migration;
\item
  production plugin/runtime changes;
\item
  methodology-neutrality verdict;
\item
  STRIDE/LINDDUN comparison or implementation.
\end{itemize}

\subsection{Commit discipline}\label{commit-discipline}

This checkpoint should be committed \textbf{before} authoring the
facial-access Decision corpus. Later empirical phases should use
separate commits/checkpoints so that the evidence sequence remains
reconstructable:

\begin{Shaded}
\begin{Highlighting}[]
\NormalTok{checkpoint: Decision construction candidate + plan}
\NormalTok{    {-}\textgreater{} facial{-}access Decision authoring}
\NormalTok{    {-}\textgreater{} facial{-}access analysis / possible regression}
\NormalTok{    {-}\textgreater{} MR{-}0003 holdout}
\NormalTok{    {-}\textgreater{} MR{-}0004 holdout}
\NormalTok{    {-}\textgreater{} MR{-}0005 holdout}
\NormalTok{    {-}\textgreater{} Decision synthesis}
\end{Highlighting}
\end{Shaded}

The purpose is methodological traceability: later revisions should make
it possible to identify which evidence changed which part of the model.

\end{document}
